% R. Brown's translations of two letters to L. Breen, with his note and comments (cote 134-2): archive cote 134-2, site page 126 (PDF page 127).
% translation of the letter of 17/2/1975. Draft compared: none.
% AI-transcribed by Claude Opus 5.5 (claude-opus-5-5, adaptive thinking, effort high), Message Batches API, batch brown_full (msgbatch_01Md3i5yK1PxkpJn5GHfEExE), 2026-11-13; not yet checked by eye.
\scanpage{134-2/126}
\mspage{- 5 -}

One has an (n+1)-functor between (n+1)-groupoids
\[ f_{n+1} : C_{n+1} \longrightarrow D_{n+1} \]
which induces by truncation an n-functor
\[ f_n : C_n \longrightarrow D_n \, . \]
One must then have:

(A') $f_n$ is an n-equivalence if and only if the n-functor
\[ f_n^{*} : \mathit{Hom}(D_n, ((n-1)\text{-Cat})) \longrightarrow \mathit{Hom}(C_n, ((n-1)\text{-Cat})) \]
which sends the local (n-1)-systems on $D_n$ (or, equally, on $D_{n+1}$)
to the local (n-1)systems on $C_n$ , is an n-equivalence.

(B') $f_n$ is an n-equivalence if and only if for every local n-system $F$
on $D_{n+1}$ ,
\[ F : D_{n+1} \longrightarrow (n\text{-Cat}) \ , \]
the n-functor induced by $f_{n+1}$
\diagram{
\begin{tikzcd}
\mathit{Hom}(e_n^{D_{n+1}}, F) \arrow[r] \arrow[d, equal, "\text{def.}"] & \mathit{Hom}(e_n^{C_{n+1}}, f^{*}_{n+1} F) \arrow[d, equal, "\text{def.}"] \\
\Gamma_{D_{n+1}}(F) & \Gamma_{C_{n+1}}(F)
\end{tikzcd}
}
is an n-equivalence.

The construction of the cohomology of a topos in terms of integration
of champs makes no appeal at all to complexes of abelian sheaves and still
less to the technique of injective resolutions. One has the impression that
in this spirit, \textit{via} the definition (which remains to be made explicit!) of
n-champs, it is all related above all to the ``\v{C}echist'' calculations in terms
of hypercoverings. Now these last are written with the help of a small dose
of semi-simplicial algebra. I do not know if a theory of champs and of
operations on them can be written \textit{without} ever using semi-simplicial algebra.
If yes, there would be essentially three distinct approaches for constructing
the cohomology of a topos:

(a) viewpoint of complexes of sheaves, injective resolutions, derived
categories (\textit{commutative homological algebra})

(b) viewpoint \v{C}echist or semi-simplicial (\textit{homotopical algebra})

(c) viewpoint of n-champs (categorical algebra, or \textit{non-commutative homological
algebra}).

In (a) one ``resolves'' the coefficients, in (b) one resolves the base
space (or topos), and in (c) it appears one resolves neither the one nor the
other.

Very cordially,

Alexandre.

\mspage{\hand{127}}
